\subsection{Early European phylogeography of SARS-CoV-2}

Following the emergence of SARS-CoV-2 in China at the end of 2019, Europe experienced large-scale COVID-19 outbreaks in early 2020.
Here we investigate early global spread through phylodynamic analyses.
Building on prior studies that incorporate international air-travel data as a predictor of migration rates, we employ a flexible unsupervised Bayesian neural network to describe the relationship between travel volumes and migration, instead of the GLM framework used previously \citep{valenzuela2021comprehensive}.
While we did not perform a formal model comparison between BELLA and the GLM, both approaches yielded similar phylogeographic reconstructions and nearly identical phylogenetic likelihoods (Fig.~\ref{fig:eucovid:likelihood}), indicating comparable overall fit to the data. However, BELLA provides a more interpretable mapping between predictors and migration rates by revealing a saturation effect at high predictor values, which is consistent with limiting processes such as finite connectivity, constrained movement opportunities, or diminishing marginal effects of the predictor once a system becomes saturated.

Based on the BELLA output, our analysis employs a phylogenetic particle filtering algorithm to sample epidemic trajectories from the posterior distribution of the population history \cite{vaughan2019estimating, vaughan2024estimates}, thereby extending earlier work \citep{nadeau2021origin}. 
The inferred transmission history indicates that initial introductions into Europe occurred around January 2020 (Fig.~\ref{fig:eucovid:summary}c), earlier than suggested by confirmed case reports.
This timing aligns both with previous phylodynamic reconstructions \citep{valenzuela2021comprehensive} and with retrospective serological evidence suggesting that SARS-CoV-2 may have been circulating in parts of Europe as early as December 2019 \citep{carrat2021evidence, amendola2021evidence}. 
Comparing the migration fluxes inferred by BELLA and the GLM reveals that, although the two models exhibit comparable measures of central tendency, they differ in their representation of uncertainty (Fig.~\ref{fig:eucovid:migrations:time_bin_2}). The GLM produces more constrained and sharply peaked distributions, whereas BELLA yields wider and flatter distributions, providing a more nuanced characterization of migration patterns and better capturing the uncertainty typical of early-stage phylodynamic epidemiological reconstructions.